% =========================================
% Chapter X — Methodology
% =========================================
\chapter{Methodology}

\section{Overview}

The methodology adopted in this thesis follows a structured three stage workflow for 
designing and evaluating AES architectures of increasing optimization complexity. The 
workflow begins with the development of a conventional AES design, progresses to an 
architecturally optimized implementation based on reference literature, and concludes with 
a power aware extension achieved through clock gating. Each stage is followed by validation 
and comparative analysis to ensure correctness and measurable improvement.

\section{Methodology Flow}

\begin{figure}[!htpb]
\centering

% ---------- LEFT SIDE (PART 1) ----------
\begin{minipage}[t]{0.48\textwidth}
\centering
\begin{tikzpicture}[
    node distance = 1.0cm,
    every node/.style = {font=\small},
    box/.style = {
        rectangle,
        rounded corners,
        draw,
        align=center,
        minimum width=5.8cm,
        minimum height=1.2cm
    },
    arrow/.style = {
        -{Latex[length=3mm]},
        thick
    }
]

\node[box] (step1) {Implementation of Conventional AES \\ (AES-128, AES-192, AES-256)};
\node[box, below=of step1] (verify1) {Functional Verification};
\node[box, below=of verify1] (step2) {Implementation of Optimized Architecture};

\draw[arrow] (step1) -- (verify1);
\draw[arrow] (verify1) -- (step2);

\end{tikzpicture}
\caption*{Baseline AES and Optimized Architecture}
\end{minipage}
\hfill
% ---------- RIGHT SIDE (PART 2) ----------
\begin{minipage}[t]{0.48\textwidth}
\centering
\begin{tikzpicture}[
    node distance = 1.0cm,
    every node/.style = {font=\small},
    box/.style = {
        rectangle,
        rounded corners,
        draw,
        align=center,
        minimum width=5.8cm,
        minimum height=1.2cm
    },
    arrow/.style = {
        -{Latex[length=3mm]},
        thick
    }
]

\node[box] (compare1) {Comparative Analysis \\ (Conventional vs. Optimized)};
\node[box, below=of compare1] (step3) {Implementation of Clock Gated Architecture};
\node[box, below=of step3] (compare2) {Final Comparative Analysis \\ (All Architectures)};

\draw[arrow] (compare1) -- (step3);
\draw[arrow] (step3) -- (compare2);

\end{tikzpicture}
\caption*{Clock Gating and Final Comparison}
\end{minipage}

\caption{Side-by-side methodology flowcharts for AES architecture design}
\label{fig:methodology_flow_sidebyside}

\end{figure}


\section{Extended Methodology for Pipeline and FSM Architectures}

The methodology described above applies to the conventional AES baseline. The present work
extends the flow to cover two additional architectures and two optimization phases.

\begin{enumerate}
    \item \textbf{Pipeline architecture implementation and verification.}
          The full-pipeline AES architecture is verified using the
          same NIST test vectors as the conventional design, using a dual-mode testbench
          controlled by the pipeline enable parameter.
          Vivado synthesis is run to capture the AES-128 pipeline baseline of 38,766 LUTs
          and 2,763 flip-flops before any RTL modifications.

    \item \textbf{FSM architecture implementation and verification.}
          The FSM AES architecture is verified in the same testbench with pipeline mode disabled.
          The AES-128 FSM baseline of
          6,794 LUTs and 412 flip-flops is captured before optimization begins.

    \item \textbf{Phase~1 --- xtime replacement and table removal.}
          The MixColumns and InvMixColumns modules are modified to replace
          EXP3/LN3 GF multiplication with \textit{xtime} shift-and-XOR arithmetic~\cite{wolkerstorfer_xtime_2002}.
          The 512-entry EXP3 and LN3 lookup-table arrays are removed from the design.
          All modules are re-verified after modification to confirm functional equivalence.

    \item \textbf{Phase~2 --- Clock enable gating on FSM registers.}
          CE guards are applied to the packed register structs in the three FSM modules.
          The cipher and inverse cipher guards freeze the register when
          \texttt{r.state == 0 \&\& Enable == 0 \&\& r.ready == 0}.
          The key expansion guard uses \texttt{r.state == 0 \&\& Enable == 0}.
          Synthesis is re-run to quantify switching power reduction.

    \item \textbf{Comparative analysis.}
          Synthesis results for the pipeline and FSM architectures are compared at baseline,
          after Phase~1, and after Phase~2 to quantify the contribution of each optimization
          stage to LUT reduction and dynamic power reduction.
\end{enumerate}

\section{Step wise Methodology Description}

\subsection*{1. Implementation of Conventional AES Architectures}

The first stage consists of designing complete AES one hundred twenty eight, AES one hundred ninety two, and 
AES two hundred fifty six architectures following the official specification. This establishes a 
baseline hardware model built from standard transformations including SubBytes, ShiftRows, 
MixColumns, and AddRoundKey, along with their inverse counterparts for decryption. The 
purpose of this stage is to provide a correct and reliable reference implementation.

\subsection*{2. Functional Verification of Conventional AES}

After implementation, the architecture is verified using standard AES test vectors and simulation 
outputs. This validation step confirms functional correctness for all three key sizes and ensures 
that subsequent optimization stages build upon a stable baseline.

\subsection*{3. Implementation of Optimized Architecture}

The second stage introduces architectural improvements based on the referenced optimized 
AES design. Key enhancements include the use of a unified S Box and inverse S Box 
structure, a combined MixColumns AddRoundKey unit, and a word serial datapath. These 
modifications focus on reducing redundant computation and improving hardware reuse while 
retaining functional equivalence with the conventional design.

\subsection*{4. Comparative Analysis of Conventional and Optimized Designs}

Both the baseline and optimized architectures are synthesized under identical conditions. 
Metrics such as logic utilization, register count, critical path, and estimated switching activity 
are evaluated. This comparison quantifies the improvements achieved through structural 
optimization and provides insight into the efficiency of the modified datapath.

\subsection*{5. Implementation of Clock Gated Architecture}

The final design stage applies clock gating to selected components of the optimized 
architecture. This extension targets the reduction of dynamic power by limiting unnecessary 
clock switching in units that are not active during every cycle. The purpose is to enhance the 
energy efficiency of the design without altering core functionality.

\subsection*{6. Comparative Analysis of All Three Architectures}

A final comparative evaluation is performed across the conventional, optimized, and clock
gated designs. The analysis examines changes in area, timing, and estimated power to
demonstrate the cumulative benefit of architectural refinement and power aware enhancement.
This comprehensive assessment validates the methodology and highlights the impact of each
stage in the design process. This includes both the pipeline and FSM architecture variants and their
respective optimization trajectories across Phase~1 and Phase~2.

